 \documentclass[12pt]{article} \usepackage[utf8]{inputenc} \usepackage[margin=1in]{geometry}

\begin{document}
\sloppy
\emergencystretch=1em

\begin{center} {\Large \textbf{INCIDENT REPORT}} \end{center} \vspace{0.5cm}

\section*{Incident Overview}

\noindent\textbf{Date of Incident:} June 12 at 11:26 AM (2026)
\vspace{0.5cm}

\noindent The PPLN crystal for the FMCW Maritime Lidar System was snapped roughly in half during an incident in the lab.
On Wednesday morning, June 10, upon visual inspection of the Lidar system Tyler Celotta and Dr. Roddewig determined
that the crystal housing was not square with both the table nor the laser setup. To fix this, the machine shop manager
came to help troubleshoot on Thursday, June 11. It was determined that the face plate connecting the housing to the
three axis adjustment needed to be faced to fit properly. This action required the removal of the housing.
On Friday morning, Tyler reassembled the setup with the newly machined part, this involved fastening the housing
to the plate using two screws, each placed in parallel to the table. After the first screw (on the left) was put in,
he held up the other end to prepare for the second screw. While reaching for the second screw, Tyler tipped the housing
toward his hand, causing the crystal tip to slide and ultimately touch his hand. In an effort to move his hand away
from the crystal, he moved his hand down slightly, which caused the crystal to flex and then break. After this,
Tyler inspected the area looking for both halves of the crystal as well as any loose fragments. He placed one broken
half on optical cloth and left the other half in the housing. Tyler then contacted Dr. Roddewig to report the incident
and discuss next steps. No personal injury occurred.

\vspace{0.5cm}
\noindent \textbf{Reported By:} Tyler Celotta \hfill \textbf{Date:} June 15, 2026

\end{document}